\subsection*{Appendix: passages in which 農 \textit{nəwṅ} `agriculture' is a rhyme word}

%Western Han according to 于安瀾 
%Yu Anlan
%冬部平 Dong Rhyme Group, Level Tone

%Western Han according to 羅 Luo Changpei and 周祖謨 Zhou Zumo
%冬部平 Dong Rhyme Group, Level Tone

（漢文, 卷25-12a）

誡子 `Admonition to My Son' by 東方朔
Dongfang Shuo (154—93 BCE）with rhyme annotation by 于安瀾 
Yu Anlan (1902—1999). The first stanza is what is relevant for us. 


\begin{longtable}{p{0.30\textwidth} p{0.60\textwidth}}
%\textbf{卷25-12a} & \\

% Line 1
明者處世, & The wise live in the world,\\
莫尚于A中 (\textit{ṭiuwṅ})。%【冬平】合韻 1)。 
& nothing surpasses balance.\\
%(Dong Ping rhyme)

% Line 2
優哉游哉, & Leisurely and carefree,\\
與道相A從 (\textit{dziəwṅ})。%【東平】。 
& following the Way.\\
%(Dong Ping rhyme)

% Line 3
首陽爲拙, & Shouyang is clumsy,\\
柳惠%2)
爲A工(\textit{kuwṅ})。 %【東平】 
& Liuhui is skilled.\\
%(Dong Ping rhyme)

% Line 4
飽食安步, & Eating fully and walking leisurely,\\
以仕代A農 (\textit{nəwṅ})。%【冬平】合韻 3) 
& serving replaces farming.\\
%(Dong Ping rhyme)

% Line 5
依隱玩世, & Secluded and playing with the world,\\
詭時不A逢 (\textit{biəwṅ})。%【東平】 
& evading unfavorable times.\\
%(Dong Ping rhyme)

%% Line 6
%是故才盡者身危, & Thus, the talented face danger;\\
%好名者得B華。%【歌平】 
%& fame-seekers gain superficial beauty.\\
%(Ge Ping rhyme)

%% Line 7
%有群者累生, & Those in groups accumulate life;\\
%孤貴者失B和。%【歌平】 
%& the solitary noble lose harmony.\\
%%(Ge Ping rhyme)

%% Line 8
%遺餘者不匱, & Those who leave something behind remain abundant;\\
%自盡者無B多。%【歌平】。 
%& those who exhaust themselves have little.\\
%%(Ge Ping rhyme)

%% Line 9
%聖人之道, & The way of the sage\\
%一龍一B蛇。%【歌平】。 
%& is like a dragon and a snake.\\
%(Ge Ping rhyme)

% Line 10
%形見神藏, & Form appears while spirit is hidden,\\
%與物變B化。%【歌平】。 
%& changing with things.\\
%%(Ge Ping rhyme)

%% Line 11
%隨時之宜, & Adapt to the times;\\
%無有常B家。%【歌平】4)。
%& have no fixed home.\\
%%(Ge Ping rhyme)

\end{longtable}


There is even actually a variant reading of the rhyme word that I am looking at. 

It looks like Zhou Zumo doesn't include 中 in the first rhyme, but otherwise they are in agreement. 

%1) 嚴注：《漢書・東方朔傳》贊作``上容''，應劭曰：容身遠害也。《御覽》作``忠''。\\
%Yan's note: The Han Shu biography of Dongfang Shuo records ``shangrong.'' Ying Shao says: ``to keep oneself safe and avoid harm.'' Yulan records ``zhong.''\\

%2) 嚴注：《漢書》作``柱下''，《御覽》作``柳下''。\\
%Yan's note: The Han Shu records ``zhuxia,'' Yulan records ``liuxia.''\\

%3) 嚴注：《漢書》作``易農''。\\
%Yan's note: The Han Shu records ``yi nong'' (to change farming).\\

%4) 嚴注：《蓺文類聚》二十三、《御覽》四百五十九。\\
%Yan's note: Cited in Yiwen Leiju vol. 23, Yulan vol. 459.\\




于安瀾 (Yu Anlan)

楊雄 (Yang Xiong)

\begin{longtable}{p{0.39\textwidth} p{0.80\textwidth}}
%世稱東方生之盛也， & The world praises the greatness of Dongfang Sheng.\\

%言不純師，行不純表， & His words were not purely those of a master, nor his actions purely exemplary.\\

%其流風、遺書，蔑如也。 & His influence and the texts he left behind are insignificant.\\

%或曰：``隱者也。'' & Someone said: "He was a recluse."\\

%曰：``昔之隱者，吾聞其語矣， & He replied: "The recluses of old—I have heard their words,\\

%又聞其行矣。'' & and I have also heard of their deeds."\\

%或曰：``隱道多端。'' & Someone said: "There are many ways to be a recluse."\\

%曰：``固也！ & He replied: "Indeed!\\

%聖言聖行，不逢其時，聖人隱也。 & If sage words and sage deeds do not meet their time, the sage remains hidden.\\

%賢言賢行，不逢其時，賢者隱也。 & If worthy words and worthy deeds do not meet their time, the worthy remain hidden.\\

%談言談行，而不逢其時，談者隱也。 & If idle words and idle deeds do not meet their time, the talker remains hidden.\\

%昔者箕子之漆其身也， & In ancient times, Jizi lacquered his body,\\

%狂接輿之被其發也， & Mad Jieyu let his hair fall loose,\\

%欲去而恐罹害者也。 & because they wished to withdraw yet feared harm.\\

%箕子之《洪範》，接輿之歌鳳也哉！'' & Jizi's "Hongfan" and Jieyu's "Song of the Phoenix!"\\

%\\[-6pt] % (Optional extra spacing between paragraphs)

%或問：``東方生名過實者，何也？'' & Someone asked: "Why is Dongfang Sheng's reputation greater than his reality?"\\

%曰：``應諧，不窮， & He replied: "He was witty, not profound.\\

%正諫，穢德。 & He gave upright remonstrance yet had sullied virtue.\\

%應諧似優，不窮似哲， & His wit seemed like jest, his lack of depth seemed like wisdom.\\

%正諫似直，穢德似隱。'' & His upright remonstrance seemed direct, his sullied virtue seemed like reclusion."\\

%\\[-6pt]

``請問名。'' & "What was his name?"\\

曰：``詼達。'' & He replied: "Huida."\\

``惡比？'' & "To whom can he be compared?"\\

曰：``非夷、齊，而是柳下惠， & He replied: "Not to Yi and Qi, but to Liuxia Hui,\\

戒其子以尚A容 (\textit{yiəwṅ}),
%【東平】， 
& who admonished his son to uphold grace. %(Dong Ping Rhyme)
\\

首陽為拙, 
%【東平】， 
& Shouyang was clumsy, 
\\

柱下為A工(\textit{kuwṅ}),
%【東平】， 
& Zhuxia was skilled. %(Dong Ping Rhyme)
\\

飽食安坐, 
%【冬平】合韻， 
& Eating fully and sitting at ease, 
\\

以仕易A農 (\textit{nəwṅ}),
%【冬平】合韻， 
&  replacing farming with official service. %(Dong Ping Rhyme)
\\

依隱玩世,
& Living in seclusion and playing with the world, 
\\


詭時不A逢(\textit{biəwṅ}),
%【東平】， 
&  missing the right time. %(Dong Ping Rhyme)
\\

其滑稽之A雄(\textit{ḫiuwṅ})
%【冬平】合韻 
乎！'' & Truly, he was a master of humor! %(Dong Ping Rhyme)
\\

%\\[-6pt]

%或問：``柳下惠非朝隱者與？'' & Someone asked: "Wasn't Liuxia Hui a recluse in the court?"\\

%曰：``君子謂之不恭。 & He replied: "The noble-minded say he was not respectful.\\

%古者高餓顯, 下祿隱。'' & In ancient times, the highly virtuous remained poor but renowned, while the lower officials withdrew into obscurity."\\

\end{longtable}







Eastern Han according to 于安瀾 Yu Anlan 
冬部平 Dong Rhyme Group, Level Tone

風陰淋農任心音潛參 

Hou Hanwen: Ban Gu, "Dou Jiangjun Beizheng Song"; Xing Nong 

Hou Hanwen: Cui Yin, "Dongxun Song"）

Eastern Han according to 羅 Luo Changpei and 周祖謨 Zhou Zumo
冬部平 Dong Rhyme Group, Level Tone


Hou Hanwen: Cui Yin, "Dongxun Song"）


竇將軍北征頌 ``Eulogy for General Dou's Northern Expedition''
by 班固 Ban Gu (32–92 CE) 
with the rhyme annotations of 
于安瀾 
Yu Anlan
(1902—1999)


全後漢文韻讀 101,  %班固于安瀾羅周周祖謨
卷26-2a

The poem is a bit long and mostly irrelevant. The section relevant here is the end of the poem.  

\begin{longtable}{p{0.40\textwidth} p{0.75\textwidth}}
%然而唱呼鬱憤， & However, they called out in repressed indignation,\\
%未逞厥A願。%【元去2】 
%& their aspirations unfulfilled.\\

%甘平原之酣A戰, %【元去2】， 
%& They relished fierce battles on the plains,\\
%矜訊捷之累筭。%【元去2】 
%& boasting of successive victories.\\

%何則？ & Why so?\\

%上將崇至仁， & The supreme general upheld utmost benevolence,\\
%行凱B易,%【錫入】合韻， 
%& making his triumphs easy.\\

%弘濃恩， & He bestowed profound kindness,\\
%降溫B澤。%【鐸入】。 
%& spreading his gentle grace.\\

%同庖廚之珍饌， & He shared the delicacies of his kitchen\\
%分裂室之纖B帛。%【鐸入】
%& and the finest silks of his chambers.\\

%勞不御輿， & Despite hardships, he did not ride in carriages;\\
%寒不施B襗。%【鐸入】。 
%& despite the cold, he did not wear extra layers.\\

%行無偏勤， & While marching, no one was overburdened;\\
%止無兼B役。%【錫入】合韻。 
%& while resting, no one was doubly tasked.\\

%悂蒙識而愎戾C順,%【真平】， 
%& The ignorant became aware, the stubborn grew obedient,\\
%貳者異而懦夫C奮。%【真平】。
%& the hesitant gained resolve, and the cowardly took courage.\\

%遂踰涿邪， & Thus, they crossed Zhuo\\
%跨祁D連,%【元平】， 
%& and traversed Qilian.\\

%籍□庭蹈就, D疆
%【陽平】合韻 
%獦崝嵮,  & They marched over enemy courtyards and trod upon rugged terrain,\\
%轔幽D山,%【元平】， 
%& rolling over secluded mountains.\\

%\symbol{"27F28}凶河， & They crossed the treacherous river,\\
%臨安候， & reached the An Garrison,\\
%軼焉居與虞D衍。%【元平】。
%& surpassing the dwellings of Yu Yan.\\

%顧衞霍之遺迹， & Looking upon the remnants of Wei and Huo,\\
%賊伊袟之所邈。 & they saw the distant domains of the defeated Yi Zhi.\\

%師橫騖而庶御， & The army advanced unhindered,\\
%士怫㥜以爭E先。%【真平】。 
%& officers and soldiers competed to be at the front.\\

%囘萬里而風騰， & Returning from thousands of miles away,\\
%劉殘寇于沂E垠。%【真平】1)。 
%& they routed the remaining enemies at the Yi border.\\

%糧不賦而師贍， & Supplies were sufficient without new levies;\\
%役不重而備E軍。%【真平】。 
%& military preparations were complete without excessive burdens.\\

%行戎醜以禮教， & The army treated captured enemies with propriety and discipline,\\
%炘鴻校而昭E仁。%【真平】。
%& demonstrating their righteousness.\\

%文武炳其並隆， & Both civil and military affairs flourished together,\\
%威德兼而兩E信%【真平】。 
%& combining authority and virtue with trust.\\

%清乾鈞之攸冒， & The purity of governance spread everywhere,\\
%拓畿略之所E順%【真平】。 
%& expanding and stabilizing the frontier.\\

%橐弓鏃而戢戈， & They sheathed their bows and arrows,\\
%囘雙麾以東E運。%【真平】。 
%& lowered their banners, and turned eastward.\\

%于是封燕然以降高， & Thus, they erected a stele at Yanran,\\
%襢廣鞬以弘F曠。%【陽去】。 
%& extending their dominion far and wide.\\

%銘靈陶以勒崇， & They inscribed their deeds in sacred stone,\\
%欽皇祇之祐F貺。%【陽去】。 
%& paying homage to the divine blessings.\\

宣惠氣， & They spread auspicious harmony,\\
盪殘G風(\textit{piuwṅ})。%【侵平】。 
& sweeping away the remnants of chaos.\\

軻泰幽嘉， & Winter settled in tranquility,\\
凝G陰 (\textit{'im})%【侵平】
飛雪。 & bringing snowfall over the frozen lands.\\

讓庶其雨， & The generous rains drenched the barren land,\\
洒G淋(\textit{lim}) %【侵平】
榛枯一握興。 & rejuvenating vegetation.\\

以仕代A農 (\textit{nəwṅ}),%【冬平】合韻， 
& Blessed plants began to grow,\\

土膏含養， & The rich earth nurtured life,\\
四行分G任(\textit{ñimH})。%【侵平】 
& sustaining all corners of the land.\\

于是三軍稱曰： & Then, the three armies proclaimed:\\

亹亹將軍， & Diligent and valiant, our general\\
克廣德G心(\textit{sim})。%【侵平】 
& broadens his virtuous heart.\\

光光神武， & Radiant and mighty,\\
弘昭德G音(\textit{'im})。%【侵平】。 
& he spreads his illustrious virtue.\\

超兮首天G潛(\textit{dziem}),%【侵平】， 
& Sublime as he rises to the heavens,\\
眇兮與神G參(\textit{śim})。%【侵平】2)。 
& vast as he communes with the divine.\\

\end{longtable}

%注釋 (Annotations)
%1) 原作``根''，《古文苑》作``垠''，于氏依《古文苑》改之。\\
%Original text has "根"; "Gu Wen Yuan" records it as "垠"; Yu corrected it based on "Gu Wen Yuan."\\

%2) 嚴注：《古文苑》，《蓺文類聚》六十九。\\
%Yan's note: Found in "Gu Wen Yuan" and "Yi Wen Lei Ju," volume 69.\\

全後漢文韻讀 194. 
"Complete Rhymed Readings of Later Han Texts, No. 194

10. 東巡頌 崔駰
"Ode on the Eastern Tour" by Cui Yin (d. 92 CE)

卷44-6a
"Volume 44, Folio 6a", with rhymes annotated by 
羅周



 
\begin{longtable}{p{0.40\textwidth} p{0.75\textwidth}}

%\multicolumn{2}{l}{0.東巡頌　崔駰}\\
%\multicolumn{2}{l}{卷44-6a}\\

%伊漢中興三葉， &
%Behold the Han, thrice renewed since its restoration;\\

%於皇維烈， &
%How majestic its blazing virtue,\\

%允迪厥倫。 &
%fully enacting the proper order.\\

%纘王命， &
%He continues the royal mandate,\\

%徹漢勳【真平】， &
%carrying forward Han’s meritorious deeds,\\

%矩坤度以範物， &
%Measuring by Earth’s pattern to guide all beings,\\

%規乾則以陶鈞【真平】。 &
%modeling after Heaven’s design to shape and mold them.\\

%于是考上帝以質中， &
%Thus, he communes with the Supreme Deity in pure sincerity,\\

%總列宿于北辰【真平】。 &
%gathering the myriad constellations around the Northern Star.\\

%開太微， &
%He opens the Grand Tenuity,\\

%敞紫庭【庚平】合韻， &
%and clears the Purple Court.\\

%延儒林， &
%He invites the community of scholars,\\

%以咨詢岱��之事。 &
%to consult on matters concerning Mount Tai.\\

%于是典司耆耇， &
%Then, the stewards and venerable elders,\\

%載華抱實， &
%bearing abundant wisdom,\\

%逌爾而造曰： &
%came forth, declaring:\\

%盛乎大漢， &
%``How splendid is the Great Han!\\

%既重雍而襲熙， &
%It renews the grand era and inherits its radiance;\\

%世增其德， &
%each generation amplifies its virtue,\\

%唯斯嶽禮， &
%but the rites of this sacred peak\\

%久而不修。 &
%have long gone unperformed.\\

%此神人之所慶幸， &
%Gods and mortals alike rejoice in this,\\

%海內之所想思。 &
%and all within the seas yearn to behold it.\\

%頌有喬山之征， &
%In hymns, there are records of ascending lofty mountains;\\

%典有徂嶽之巡【真平】。 &
%in the canons, accounts of touring Mount Cu.\\

%時邁其邦， &
%Now is the time to journey through the realm,\\

%民斯攸勤【真平】， &
%and let the people devote themselves with zeal—\\

%不亦宜哉！ &
%is this not most fitting?\\

%乃命太僕， &
%Thus he orders the Grand Equerry,\\

%訓六騶【魚平】， &
%to train the six teams of horses,\\

%閑路馬， &
%calm the road‐steeds,\\

%戒師徒【魚平】。 &
%and prepare the ranks.\\

%于是乘輿登天靈之威路， &
%So, the imperial carriage traverses the august path of the Heavenly Spirits,\\

%駕太一之象車【魚平】， &
%harnessing the chariot symbolic of the Great One,\\

%升九龍之華旗， &
%raising the Nine‐Dragon banners in resplendence,\\

%建翠霓之旌旄【宵平】合韻。 &
%and unfurling the rainbow‐green standards.\\

三軍霆激， &
The three armies roar like thunder,\\

羽騎火烈， &
and the winged cavalry flares with fiery might.\\

天動雷E震 (\textit{cinH}),
%【真平】, 
&
Heaven quakes with pealing thunder,\\

隱隱轔E轔 (\textit{lin}),
%【真平】。
&
its rumble rolling on and on.\\

躬東作之上務, &
He personally oversees the foremost labors of the East,\\

始八正于南F行 (\textit{ḫaṅ})。
%【陽平】。 
&
and commences the ``Eight Rectitudes'' with a journey to the South.\\

裒胡耇之元老， &
He gathers the venerable elders, even among the Hu tribes,\\

賞孝行之畯F農 (\textit{nəwṅ})。
%【冬平】合韻 1)。 
&
and rewards the exemplary farmers who uphold filial devotion.\\

\end{longtable}



According to 周祖謨 Zhou Zumo in the Jin dynasty 
宵部平 Xiao Rhyme Group, Level Tone


曹農 晉文.束皙.勸農賦、茅農 晉文.潘岳.藉田賦

Cao Nong – Jin Literature: Shu Xi, "Quan Nong Fu"; Mao Nong – Jin Literature: Pan Yue, "Jie Tian Fu"


全晉文韻讀 341. 束皙 周祖謨
Complete Rhymed Readings of Jin Literature, No. 341: Shu Xi, with annotations by Zhou Zumo.

4. 勸農賦 Rhapsody on Encouraging Agriculture by 束皙 Shu Xi
（262-300 CE）

卷87-2
Volume 87, Folio 2


\begin{longtable}{p{0.40\textwidth} p{0.75\textwidth}}

%\multicolumn{2}{l}{4. 勸農賦　束皙}\\
%\multicolumn{2}{l}{卷87-2}\\

惟百里之置吏， &
For the officials appointed in each hundred \emph{li},\\

各區別而異A曹。(\textit{dzaw})
%【宵平】。 
&
they are separately classified into distinct bureaus.\\

考治民之賤職， &
Examining the lowliest role of governing the people,\\

美莫當乎勸A農。(\textit{nəwṅ})
%【宵平】
&
none is more praiseworthy than exhorting them to agriculture.\\

%專一里之權， &
%They wield authority over a single district,\\

%擅百家之勢。 &
%and command the strength of a hundred households.\\

%及至青幡禁乎游惰， &
%When the green banners forbid idleness and wandering,\\

%田賦度乎頃畝【侯上】。 &
%the land tax is measured by \emph{qing} and \emph{mu}.\\

%與奪在己， &
%Giving or withholding rests in their own hands,\\

%良薄浹口【侯上】。 &
%so favor or neglect stirs the entire populace.\\

%受饒在於肥脯， &
%They take their fill in luscious meats,\\

%得力在於美酒【侯上】。 &
%and derive their vigor from fine wines.\\

%若場功畢， &
%When the harvest tasks are complete,\\

%租輸至， &
%the dues and grain levies arrive,\\

%錄社長， &
%they record village leaders,\\

%召閭師。 &
%and summon neighborhood heads.\\

%條牒所領， &
%Then, by written orders,\\

%注列名諱。 &
%they list each name in the register.\\

%則雞豚爭下， &
%Thereupon, chickens and pigs are rushed to slaughter,\\

%壺榼橫至。 &
%and flasks and pitchers are lavishly brought forth.\\

%遂乃定一以為十， &
%Hence they fix ``one'' to be ``ten,''\\

%拘五以為二。 &
%and constrain ``five'' to be ``two.''\\

%蓋由熱啖紆其腹， &
%It is all because of hot feasting easing their bellies,\\

%而杜康咥其胃1)。 &
%and Du Kang’s ale delighting their guts.\\

%乃有老閑舊猥， &
%Yet there are those long idle or corrupt from old habit,\\

%欺狹難覺【沃入】。 &
%their deceptions subtle and hard to detect.\\

%時雖被考， &
%Even when subjected to scrutiny,\\

%不過校督【沃入】。 &
%they are merely checked or supervised.\\

%歌對囹圄， &
%They sing in the face of prison,\\

%笑向桎梏【沃入】2)。 &
%and laugh though they wear fetters.\\

\end{longtable}

全晉文韻讀 356. 潘岳
Complete Rhymed Readings of Jin Literature, No. 356: Pan Yue.

19. 藉田賦 潘岳
Rhapsody on Plowing the Fields by Pan Yue (247-300CE).

周祖謨
With Commentary by Zhou Zumo.

卷91-3b
Volume 91, Folio 3b.


\begin{longtable}{p{0.40\textwidth} p{0.75\textwidth}}

\multicolumn{2}{l}{19. 藉田賦　潘岳}\\
\multicolumn{2}{l}{卷91-3b}\\

%伊晉之四年， &
%In the fourth year of the Jin,\\

%正月丁未， &
%on the day Dingwei in the first month,\\

%皇帝親率羣后 &
%the Emperor himself led the assembled nobles\\

%藉于千畝之甸，禮也。 &
%to plow at the fields of a thousand \emph{mu}, in accordance with the rites.\\

%于是乃使甸帥清畿， &
%Then he ordered the Superintendent of the Fields to cleanse the imperial domain,\\

%野廬掃路【魚去】。 &
%and for the rural huts to sweep the roads.\\

%封人壝宮， &
%The wardens built boundary walls,\\

%掌舍設枑【魚去】。 &
%and the quartermasters erected enclosure gates.\\

%青壇蔚其岳立兮， &
%A verdant altar rose high as a mountain,\\

%翠幕黕以雲布【魚去】。 &
%and a green canopy stretched out like clouds.\\

%結崇基之靈趾兮， &
%They laid a lofty foundation for this sacred site,\\

%啟四塗之廣阼【魚去】。 &
%opening four thoroughfares to its spacious forecourt.\\

%沃野墳腴， &
%Rich fields swelled with fertile soil,\\

%膏壤平砥【脂上】。 &
%the fat earth level as a whetstone.\\

%清洛濁渠， &
%Clear waters of the Luo, turbid ditches,\\

%引流激水【脂上】。 &
%direct the flow and channel the currents.\\

%遐阡繩直， &
%Far‐stretching paths are ruler‐straight,\\

%邇陌如矢【脂上】。 &
%and near lanes align like arrows.\\

%��犗服于縹軛兮， &
%Oxen yoked in azure harnesses,\\

%紺轅綴于黛耜【之上】合韻。 &
%dark shafts affixed to dusky plowshares.\\

%儼儲駕于廛左兮， &
%In neat array, the teams are readied to one side of the settlement,\\

%俟萬乘之躬履【脂上】。 &
%awaiting the Emperor himself to tread the fields.\\

%百僚先置， &
%First come the hundred officials,\\

%位以職分【真平】。 &
%each assigned rank according to duty.\\

%自上下下， &
%From highest to lowest,\\

%具惟命臣【真平】。 &
%all are subject to the imperial mandate.\\

%襲春服之萋萋兮， &
%They don the lush garments of spring,\\

%接游車之轔轔【真平】。 &
%accompanied by the rumbling of traveling carriages.\\

%微風生于輕幰， &
%A light breeze stirs beneath their silken canopies,\\

%纖埃起于朱輪【真平】。 &
%fine dust rises around the crimson wheels.\\

%森奉璋以階列， &
%They form ranks, each bearing jade emblems,\\

%望皇軒而肅震【真平】。 &
%looking upon the Emperor’s chariot in solemn awe.\\

%若湛露之晞朝陽， &
%Like dewdrops glistening in the morning sun,\\

%似眾星之拱北辰【真平】也。 &
%as if myriad stars array themselves around the northern pole.\\

%于是前驅魚麗， &
%Then the foremost procession forms like a school of fish,\\

%屬車鱗萃【脂去】。 &
%carriages joined in ranks like scales.\\

%閶闔洞啟， &
%The palace gates stand wide,\\

%參塗方駟【脂去】。 &
%and triple‐layered roads host chariots drawn by square‐arranged fours.\\

%常伯陪乘， &
%The Chief Equerry attends the imperial ride,\\

%太僕秉轡【脂去】。 &
%while the Grand Coachman holds the reins.\\

%后妃獻穜稑之種， &
%The Empress and consorts present seeds of millets and grains,\\

%司農撰播殖之器【脂去】。 &
%as the Minister of Agriculture selects the tools for sowing.\\

%挈壺掌升降之節， &
%The Timekeeper manages the timing of up and down beats,\\

%宮正設門閭之蹕【質入】合韻。 &
%and the Palace Overseer sets up the watchwords at the gates.\\

%天子乃御玉輦， &
%Then the Son of Heaven mounts the jade carriage,\\

%蔭華蓋【泰去】。 &
%beneath a canopy splendid as a blossoming tree.\\

%衝牙錚鎗， &
%Its ornamented metal bits clang sharply,\\

%綃紈綷䌨【曷入】合韻。 &
%and the silk harness straps stretch taut.\\

%金根照耀以烱晃兮， &
%Golden axles gleam in bright luster,\\

%龍驥騰驤而沛艾【泰去】。 &
%while dragon steeds bound forward in galloping force.\\

%表朱玄于離坎， &
%They display red and black standards marking south and north,\\

%飛青縞于震兌【泰去】。 &
%and hoist green‐white banners at east and west.\\

%中黃曄以發揮， &
%A vivid central yellow bursts forth,\\

%方綵紛其繁會【泰去】。 &
%as swirling colors converge in abundance.\\

%五輅鳴鑾， &
%Five state chariots ring their bells,\\

%九旗揚旆【泰去】。 &
%nine banners unfurl their pennants.\\

%瓊鈒入⿱蕊木， &
%Jade fittings adorn the blossoming carriage wood,\\

%雲罕晻藹【泰去】。 &
%while cloudlike canopies cast a dusky shade.\\

%簫管嘲哳以啾嘈兮， &
%Flutes and pipes twitter and trill,\\

%鼓鞞硡隱以砰礚【泰去】。 &
%as drums and gongs resound with booming thunder.\\

%筍簴嶷以軒翥兮， &
%Stately balustrades loom high like soaring wings,\\

%洪鐘越乎區外【泰去】。 &
%and the great bell echoes beyond all bounds.\\

%震震填填【先平】， &
%Loud and unceasing the tumult,\\

%塵騖連天【先平】， &
%dust billowing until it meets the heavens,\\

%以幸乎藉田【先平】。 &
%as they proceed to grace the ``Plowing Fields.''\\

%蟬冕熲以灼灼兮， &
%Their cicada‐shaped coronets glow radiantly,\\

%碧色肅其千千【先平】。 &
%while thousands clad in green stand in reverent ranks.\\

%似夜光之剖荆璞兮， &
%Like a luminous gem revealed from its raw jade,\\

%若茂松之依山巔【先平】也。 &
%or a thriving pine perched upon the mountain crest.\\

%于是我皇乃降靈壇， &
%At this point, Our Sovereign descends from the sacred altar,\\

%撫御耦【侯上】。 &
%and personally takes hold of the plow handles.\\

%坻場染屨， &
%Stepping upon the soaked furrows, mud staining his shoes,\\

%洪縻在手【侯上】。 &
%the great tether firmly in hand.\\

%三推而舍， &
%With three pushes he stops,\\

%庶人終畝【侯上】。 &
%while the common folk finish the furrow.\\

%貴賤以班， &
%Noble and humble are assigned tasks,\\

%或五或九【侯上】。 &
%some in groups of five, some of nine.\\

%于斯時也。 &
%At this time,\\

%居靡都鄙， &
%the populace dwell in no big cities or remote outposts—\\

%民無華裔【祭去】。 &
%no one is labeled ``Chinese'' or ``foreigner.''\\

%長幼雜遝以交集， &
%Young and old commingle in a crowd,\\

%士女頒斌而咸戾【祭去】。 &
%men and women gather in cheerful throngs.\\

%被褐振裾， &
%Clad in coarse robes, they shake out their hems,\\

%垂髫總髻【祭去】1)。 &
%while the children with dangling locks tie their hair in topknots.\\

%躡踵側肩， &
%They press forward on tiptoe, shoulders brushing,\\

%掎裳連襼【祭去】。 &
%skirts hooking together in a bustle.\\

%黃塵為之四合兮， &
%A yellow dust rises up on every side,\\

%陽光為之潛翳【祭去】。 &
%even the sunlight grows dim and veiled.\\

%動容發音而觀者， &
%Those who watch, stirred in body and voice,\\

%莫不抃舞乎康衢， &
%cannot help but clap and dance in the joyous streets,\\

%謳吟乎聖世【祭去】。 &
%singing praises to this sacred age.\\

%情欣樂于昏作兮， &
%Their hearts delight to labor from dawn till dusk,\\

%慮盡力乎樹蓺【祭去】。 &
%striving with all might to plant and cultivate.\\

%靡誰督而常勤兮， &
%No one need oversee them, yet they work industriously,\\

%莫之課而自厲【祭去】。 &
%no one compels them, yet they spur themselves on.\\

%躬先勞以說使兮， &
%They see the Emperor’s personal toil as an encouraging model—\\

%豈嚴刑而猛制【祭去】之哉！ &
%what need is there for harsh punishments or severe edicts?\\

%有邑老田父， &
%Then, an elder from these fields,\\

%或進而稱曰： &
%stepped forward and spoke, saying:\\

%蓋損益隨時，理有常然【先平】。 &
%``In matters of gain and loss, one adapts with the times, for that is the natural order.\\

%高以下為基，民以食為天【先平】。 &
%From high to low, the foundation stands; for the people, food is heaven’s decree.\\

%正其末者端其本， &
%He who corrects the branches sets the trunk aright,\\

%善其後者慎其先【先平】。 &
%he who perfects the end carefully tends the beginning.\\

%夫九土之宜弗任， &
%If the nine kinds of terrain are not fully utilized,\\

%四人之務不壹【質入】。 &
%the tasks of all four classes cannot be unified.\\

%野有菜蔬之色， &
%Fields may yield green vegetables,\\

%朝靡代耕之秩【質入】。 &
%yet the court lacks offices to replace those who plow.\\

%無儲稸以虞災， &
%Without stored grain to guard against misfortune,\\

%徒望歲以自必【質入】。 &
%we vainly count on good harvests as though guaranteed.\\

%三季之衰，皆此物【質入】也。 &
%In the decline of the Three Ages, these very failings abounded.\\

%今聖上昧旦丕顯， &
%But now, Our Sagely Lord rises early in great splendor,\\

%夕惕若慄【質入】。 &
%and at evening remains vigilant, as if trembling.\\

%圖匱于豐， &
%He anticipates want even in times of plenty,\\

%防儉于逸【質入】。 &
%and guards against frugality’s needs in days of ease.\\

%欽哉欽哉，惟穀之䘏【質入】。 &
%How admirable—truly mindful of the grain supply.\\

%展三時之弘務， &
%Expanding the grand tasks of the three seasons,\\

%致倉廩于盈溢【質入】。 &
%he fills the granaries to overflowing.\\

%固堯湯之用心， &
%This is indeed the way of Yao and Tang,\\

%而存救之要術【質入】也。 &
%the essential art of safeguarding the people.\\

%若乃廟祧有事， &
%When events occur in ancestral temples,\\

%祝宗諏日【質入】。 &
%the praying priests select an auspicious day.\\

%簠簋普淖，則此之自實【質入】。 &
%Their sacrificial vessels are broadly replenished—these fields provide it all.\\

%縮鬯蕭茅，又于是乎出【質入】。 &
%The aromatic ale and sweet rushes likewise come forth from here.\\

%黍稷馨香，旨酒嘉栗【質入】。 &
%The millet is fragrant, the wine is delicious, and the chestnuts are fine.\\

%宜其民和年登， &
%Thus the people dwell in harmony, the harvests flourish,\\

%而神降之吉【質入】也。 &
%and the gods bestow good fortune.\\

%古人有言曰， &
%The ancients have said:\\

%聖人之德，無以加于孝乎！ &
%‘No virtue of the sage exceeds filial piety!’\\

%夫孝，天地之性， &
%Indeed, filial piety is Heaven and Earth’s own essence,\\

%人之所由靈也。 &
%the root of humanity’s spirit.\\

%昔者明王以孝治天下， &
%In former days, the enlightened kings ruled all under Heaven by this filial virtue;\\

%其或繼之者，鮮哉希矣！ &
%those who continued it were few and far between!\\

%逮我皇晉， &
%But now, under Our Imperial Jin,\\

%實光斯道【宵上】。 &
%this Way is truly brought to shining fulfillment.\\

%儀刑孚于萬國， &
%His exemplary conduct wins trust from the myriad realms,\\

%愛敬盡于祖考【宵上】。 &
%and his love and reverence fully extend to his ancestors.\\

%故躬稼以供粢盛， &
%Therefore he personally tills the fields to provide the sacrificial grains—\\

%所以致孝也！ &
%this is how he enacts filial piety!\\

%勸穡以足百姓， &
%Encouraging agriculture to enrich the populace—\\

%所以固本也。 &
%thus he secures the foundation.\\

%能本而孝， &
%To keep the foundation sound, and remain filial—\\

%盛德大業至矣哉！ &
%such is the summit of great virtue and grand endeavor.\\

%此一役也， &
%In this single undertaking,\\

%而二美具焉。 &
%two splendid ends are met.\\

%不亦遠乎， &
%Is this not far‐reaching?\\

%不亦重乎！ &
%Is this not of great import?\\

敢作頌曰。 &
I venture to compose a hymn, saying:\\

思樂甸幾， &
I think joyfully of those fertile fields,\\

薄采其L茅。
%【宵平】。 
&
and lightly gather their thatch.\\

大君戾止， &
Our great lord arrives,\\

言藉其L農。
%【宵平】。
&
to plow with the peasants in those fields.\\

其農三推， &
He plows thrice,\\

萬方以M祇。
%【支平】。
&
and myriad lands show him reverence.\\

耨我公田， &
We weed these state fields,\\

實及我M私。
%【脂平】合韻。
&
and thereby our private plots as well.\\

我簠斯盛， &
Thus our sacrificial vessels are full,\\

我簋斯M齊。
%【皆平】合韻。
&
our bowls neatly ordered.\\

我倉如陵， &
Our granaries stand like hills,\\

我庾如M坻。
%【脂平】合韻。 
&
our storehouses loom like islets.\\

%念茲在茲， &
%Holding this always in our thoughts,\\

%永言孝思【之平】。 &
%we forever cherish filial devotion.\\

%人力普存， &
%Everyone exerts their strength,\\

%祝史正辭【之平】。 &
%and the priests offer solemn words.\\

%神祇攸歆， &
%The gods partake of the fragrance,\\

%逸豫無期【之平】。 &
%and joy and ease know no end.\\

%一人有慶， &
%When the ruler attains blessing,\\

%兆民賴之【之平】2)。 &
%countless people benefit thereby.\\

\end{longtable}


%According to 于安瀾 Yu Anlan in the Song Dynasty
%東部平 (Dong Rhyme Group, Level Tone

%（風邦恭農 宋文.顏延之.陶徵士誄）\\
%\textit{(Feng, Bang, Gong, Nong – Song Literature: Yan Yanzhi, "Tao Zhengshi Lei")}

%According to 周祖謨 in the Song dynasty 


%東部平 Dong Rhyme Group, Level Tone
%（風邦恭農 宋文.顏延之.陶徵士誄）
%(Feng, Bang, Gong, Nong – Song Literature: Yan Yanzhi, "Tao Zhengshi Lei"


全宋文韻讀 85. 顏延之
Complete Rhymed Readings of Song Dynasty Literature, No. 85: Yan Yanzhi.

3. 陶徵士誄　顏延之
Eulogy for the Hermit Tao, by Yan Yanzhi (384—456 CE).

卷38-2b
Volume 38, Folio 2b.



\begin{longtable}{p{0.40\textwidth} p{0.75\textwidth}}

\multicolumn{2}{l}{3. 陶徵士誄　顏延之}\\
\multicolumn{2}{l}{卷38-2b}\\

%物尚孤生， &
%All things tend to exist in solitary form,\\

%人固介立【緝入】。 &
%and man, by nature, stands alone.\\

%豈伊時遘， &
%Is it only the era that brings them together,\\

%曷云世及【緝入】？ &
%or can the world truly unite them?\\

%嗟乎若士， &
%Alas, such a man,\\

%望古遙集【緝入】。 &
%gazing upon antiquity, gathers remote ideals.\\

%韜此洪族， &
%He conceals his grand lineage,\\

%蔑彼名級【緝入】。 &
%and disregards fame or rank.\\

%睦親之行， &
%His cultivation of kinship sprang from sincerity,\\

%至自非敦【魂平】。 &
%though he never forced it overtly.\\

%然諾之信， &
%His promise was as good as a sacred pledge,\\

%重於布言【先平】合韻 。 &
%weightier than mere spoken words.\\

%廉深簡潔， &
%His integrity was profound, his manner simple and unadorned,\\

%貞夷粹溫【魂平】。 &
%steadfast, pure, and amiably mild.\\

%和而能峻， &
%Gentle yet uncompromising,\\

%博而不繁【先平】合韻 。 &
%broad in scope yet never excessive.\\

%依世尚同， &
%Conforming to the times, he seemed at one with them;\\

%詭時則異【之去】。 &
%when rejecting the times, he stood apart.\\

%有一於此， &
%If someone tries to fuse the two,\\

%兩非默置【之去】。 &
%they end up satisfying neither, and so remain silent.\\

%豈若夫子， &
%But how unlike this gentleman,\\

%因心違事【之去】。 &
%who followed his heart and shunned worldly affairs.\\

%畏榮好古， &
%He feared false honors, loved the ancient ways,\\

%薄身厚志【之去】。 &
%held his own person lightly but his aspirations in high regard.\\

世霸虛禮， &
Amid an age of hollow decorum,\\

州壤推風【東平】。 &
regional domains feigned virtue.\\

孝惟義養， &
His filial piety sought only righteous support,\\

道必懷邦【東平】。 &
his Way ever embracing the realm.\\

人之秉彝， &
Human nature, abiding in constant principle,\\

不隘不恭【東平】。 &
is neither narrow nor servile.\\

爵同下士， &
Rank equal to the low official,\\

祿等上農【東平】。 &
stipend akin to a high peasant.\\

%度量難鈞， &
%Yet matching his breadth of mind was no easy feat;\\

%進退可限【先上】。 &
%thus, his advancement or withdrawal knew its own bounds.\\

%長卿棄官， &
%Like Changqing who cast aside his office,\\

%稚賓自免【先上】。 &
%like Zhibin who freed himself from public service.\\

%子之悟之， &
%You understood their awakenings,\\

%何悟之辯【先上】。 &
%and how truly enlightening that insight was!\\

%賦詩歸來， &
%You composed poems as you returned,\\

%高蹈獨善【先上】。 &
%ascending to lofty solitude and preserving your own virtue.\\

%亦既超曠， &
%Having transcended all confines,\\

%無適非心【侵平】。 &
%nowhere you went could be contrary to your nature.\\

%汲流舊巘， &
%You drew from the stream near your old heights,\\

%葺宇家林【侵平】。 &
%repaired your house among the family groves.\\

%晨煙暮藹， &
%Morning mists and evening haze,\\

%春煦秋陰【侵平】。 &
%spring warmth and autumn gloom—\\

%陳書輟卷， &
%you would set your books aside,\\

%置酒絃琴【侵平】。 &
%to pour wine and strum the zither.\\

%居備勤儉， &
%Your dwelling practiced diligence and thrift,\\

%躬兼貧病【庚去】。 &
%your own person enduring both poverty and illness.\\

%人否其憂， &
%Others might find that cause for sorrow,\\

%子然其命【庚去】。 &
%yet you accepted your fate as it was.\\

%隱約就閑， &
%Withdrawn and quiet, you embraced a secluded life,\\

%遷延辭聘【庚去】。 &
%lingering in reclusion, declining all summons.\\

%非直也明， &
%Not only in outward candor,\\

%是惟道性【庚去】。 &
%but truly in the nature of your Way.\\

%糾纆斡流， &
%Binding ropes can reverse a current,\\

%冥漠報施【支去】。 &
%yet in the vast unseen, deeds still find repayment.\\

%孰云與仁， &
%Who claims to share in benevolence?\\

%實疑明智【支去】。 &
%One might well question even clear wisdom.\\

%謂天蓋高， &
%We say Heaven is high,\\

%胡諐斯義。 &
%but how can such duty be evaded?\\

%履信曷憑， &
%If we tread on trust, where do we rely,\\

%思順何寘【支去】？ &
%and where place our wish to comply with Fate?\\

%年在中身， &
%At mid‐life,\\

%疚維痁疾【質入】。 &
%you were already plagued by sickness.\\

%視死如歸， &
%You looked upon death as a return home,\\

%臨凶若吉【質入】。 &
%confronting disaster as though it were good fortune.\\

%藥劑弗嘗， &
%You never tasted medicinal brews,\\

%禱祀非恤【質入】。 &
%nor troubled yourself with prayers or sacrifices.\\

%傃幽告終， &
%You fixed your course toward the darkness beyond, ending your days,\\

%懷和長畢【質入】。 &
%embracing peace to the very last.\\

%嗚呼哀哉！ &
%Alas! Our grief is profound.\\

%敬述靖節， &
%I respectfully recount your tranquil integrity,\\

%式尊遺占【鹽去】。 &
%thereby honoring the signs you have left behind.\\

%存不願豐， &
%In life, you never sought abundance,\\

%沒無求贍【鹽去】。 &
%and in death, you required no lavish provisions.\\

%省訃卻賻， &
%Hearing the sad news, you refused funeral gifts;\\

%輕哀薄斂【鹽去】。 &
%your mourning was spare, your burial modest.\\

%遭壤以穿， &
%The earth was dug for your resting place,\\

%旋葬而窆【鹽去】。 &
%and you were soon interred therein.\\

%嗚呼哀哉！ &
%Alas! How grievous it is.\\

%深心追往， &
%With a pensive heart, I trace memories past,\\

%遠情逐化【歌去】。 &
%my distant thoughts follow your transformation.\\

%自爾介居， &
%Ever since you dwelt alone,\\

%及我多暇【歌去】。 &
%and I found ample time—\\

%伊好之洽， &
%this congenial bond between us,\\

%接閻鄰舍【歌去】。 &
%joining neighbors and friends—\\

%宵盤晝憩， &
%we chatted late into the night, rested in the day,\\

%非舟非駕【歌去】。 &
%lacking neither boat nor carriage.\\

%念昔宴私， &
%Recalling those private feasts before,\\

%舉觴相誨【咍去】。 &
%we raised our cups, admonishing and guiding each other.\\

%獨正者危， &
%One who stands purely upright meets peril,\\

%至方則礙【咍去】。 &
%the perfectly square block will not fit all corners.\\

%哲人卷舒， &
%The sage may furl or unfold like a scroll,\\

%布在前載【咍去】。 &
%his words spread across the records of old.\\

%取鑒不遠， &
%No need to seek mirrors in distant times—\\

%吾規子佩【咍去】。 &
%I’ve taken your example to heart, that I might wear it like a token.\\

%爾實愀然， &
%You grew solemn in spirit,\\

%中言而發【沒入】。 &
%speaking from the depths of your being.\\

%違眾速尤， &
%You parted from the crowd and quickly drew censure,\\

%迕風先蹷【沒入】。 &
%meeting harsh gales sooner than most would stumble.\\

%身才非實， &
%Your gifts were not robust,\\

%榮聲有歇【沒入】。 &
%your acclaim not lasting.\\

%叡音永矣， &
%Your wise voice, lost forever—\\

%誰箴余闕【沒入】。 &
%who now shall admonish my faults?\\

%嗚呼哀哉， &
%Alas! Our sorrow is immense.\\

%仁焉而終， &
%He died in benevolence,\\

%智焉而斃【祭去】。 &
%his wisdom likewise perished.\\

%黔婁既沒， &
%As with Qian Lou’s passing,\\

%展禽亦逝【祭去】。 &
%like Zhan Qin’s departure—\\

%其在先生， &
%so it is with our Master,\\

%同塵往世【祭去】。 &
%gone as dust in a vanished age.\\

%旌此靖節， &
%We extol this quiet integrity,\\

%加彼康惠【祭去】。 &
%granting him well‐earned grace.\\

%嗚呼哀哉1)！ &
%Alas! Our lament knows no bounds.\\

\end{longtable}

噥 nowng doesn't occur in Suzuki's data 

膿 nowng occurs twice in pre-Qin stuff 


\begin{longtable}{p{0.30\textwidth} p{0.75\textwidth}}

\multicolumn{2}{l}{靈樞韻讀 7. 官鍼姚文田江有誥}\\
\multicolumn{2}{l}{姚文田}\\
\multicolumn{2}{l}{江有誥}\\
\multicolumn{2}{l}{1.}\\

%凡刺之要【15爻去】， &
%``In the essentials of acupuncture,\\

%官鍼最妙【15爻去】。 &
%the ‘Official Needle’ is the finest.''\\

%九鍼之宜【11麻平】， &
%``Among the nine types of needles,\\

%各有所爲【11麻平】， &
%each has its specific purpose.\\

%長、短、大、小， &
%They differ in length, thickness, and size;\\

%各有所施【11麻平】也。 &
%each is used according to need.\\

%不得其用， &
%If one fails to use them correctly,\\

%病弗能移【11麻平】。 &
%the ailment cannot be driven away.\\

疾淺鍼深， &
If the illness is shallow but the needle goes in deeply,\\

內傷良肉， &
it harms healthy tissue inside,\\

皮膚爲癰【1東平】； &
causing sores to form on the skin.\\

病深鍼淺， &
If the illness is deep but the needle is placed shallowly,\\

病氣不寫， &
the pathogenic qi is not discharged,\\

支爲大膿【1東平】。 &
and it accumulates into large abscesses.\\

%病小鍼大【5齊去】， &
%If the ailment is minor but the needle is too large,\\

%氣寫大甚， &
%the discharge of qi is excessive,\\

%疾必爲害【5齊去】； &
%and the illness inevitably worsens.\\

%病大鍼小， &
%If the ailment is severe but the needle is too small,\\

%氣不泄【5齊去】寫， &
%qi cannot be sufficiently released,\\

%亦復爲敗【5齊去】。 &
%resulting in further deterioration.\\

%失鍼之宜【11麻平】。 &
%Such is the error of misusing the needle.\\

%大者寫1)， &
%Severe cases must be drained,\\

%小者不移【11麻平】。 &
%while mild ones remain unchanged.\\

%已言其過， &
%We have spoken of these errors;\\

%請言其所施【11麻平】。 &
%let us now discuss the proper applications.\\

\end{longtable}


 靈樞韻讀 75. 刺節真邪. 　姚文田　

\begin{longtable}{p{0.30\textwidth} p{0.75\textwidth}}

\multicolumn{2}{l}{8.}\\

凡刺癰邪， &
In needling abscess‐type ailments,\\

無迎隴【1東平】， &
do not chase them beyond the ridges.\\

易俗移性， &
Change the internal environment,\\

不得膿【1東平】， &
so that no pus forms.\\

%脆道更行， &
%By altering delicate pathways,\\

%去其鄉【16庚平】， &
%the pathogen is forced from its habitat.\\

%不安處所 &
%Unable to remain in place,\\

%乃散亡【16庚平】， &
%it disperses and vanishes.\\

%諸陰陽過癰取之， &
%Should yin or yang be in excess, producing an abscess, needle it,\\

%其輸寫之。 &
%and drain the pathogenic qi at its proper points.\\

\end{longtable}



